\documentclass[11pt,a4paper]{article}

% --- Langue, polices, mise en page ------------------------------------------------------------
\usepackage[T1]{fontenc}
\usepackage[utf8]{inputenc}
\usepackage{lmodern}
\usepackage[french]{babel}
\usepackage[a4paper,margin=2.4cm,headheight=14pt]{geometry}
\usepackage{microtype}

% --- Tableaux, figures, listes ----------------------------------------------------------------
\usepackage{booktabs,tabularx,array}
\usepackage{xcolor}
\usepackage{tikz}
\usetikzlibrary{arrows.meta,positioning}
\usepackage[font=small,labelfont=bf,skip=6pt]{caption}
\usepackage{enumitem}
\usepackage{float}
\usepackage{titlesec}
\titlespacing*{\paragraph}{0pt}{5pt}{0.6em}
\setlist{itemsep=2pt,topsep=4pt}

% --- En-têtes et liens ------------------------------------------------------------------------
\usepackage{fancyhdr}
\pagestyle{fancy}
\fancyhf{}
\fancyhead[L]{\small\textsc{Atelier 2 — Tests unitaires}}
\fancyhead[R]{\small Projet \texttt{e-cheque}}
\fancyfoot[C]{\small\thepage}
\renewcommand{\headrulewidth}{0.4pt}

\definecolor{lien}{HTML}{1F4E79}
\usepackage[colorlinks=true,linkcolor=black,urlcolor=lien]{hyperref}
\hypersetup{
  pdftitle={Atelier 2 — Tests unitaires : démarche et solution},
  pdfauthor={Mehdi BELLAHSENE},
  pdfsubject={Module de communication POS (e-cheque)},
}

\newcommand{\code}[1]{\texttt{\small #1}}
\newcolumntype{L}{>{\raggedright\arraybackslash}X}
\setlength{\parindent}{0pt}
\setlength{\parskip}{4pt}

\title{\textbf{Atelier 2 — Tests unitaires}\\[4pt]
  \large Démarche et solution pour le module de communication avec la caisse}
\author{Mehdi BELLAHSENE\\[2pt]
  \small Dépôt : \href{https://github.com/mehdibellahsene/e-cheque}{github.com/mehdibellahsene/e-cheque}}
\date{18 novembre 2026}

\begin{document}

\maketitle
\thispagestyle{plain}

\begin{abstract}
\noindent Un module Rust qui interroge la caisse sur une liaison série est généré par une IA,
puis rendu testable, testé, mesuré et éprouvé par mutation. Les tests ont révélé six réponses
invalides acceptées par le code généré, toutes corrigées. Le résultat : 60~tests, une
couverture de 100\,\% des lignes et des branchements, et aucun mutant survivant sur~63.
\end{abstract}

% =============================================================================================
\section{Contexte et objectif}

Un lecteur de chèques demande à la caisse la date et le montant du chèque en cours.
L'objectif est de \emph{prouver par des tests} que le module chargé de cet échange
fonctionne, y compris lorsque la caisse répond mal. Le protocole est résumé au
tableau~\ref{tab:protocole}.

\begin{table}[H]
\centering
\caption{Protocole d'échange avec la caisse.}
\label{tab:protocole}
\begin{tabularx}{\linewidth}{@{}l l L@{}}
\toprule
Sens & Trame & Signification \\
\midrule
Lecteur $\to$ caisse & \code{GET\_CHEQUE\_DATA\textbackslash n} & Demande des données \\
Caisse $\to$ lecteur & \code{OK;2026-11-16;145.50\textbackslash n} & Réponse nominale \\
Caisse $\to$ lecteur & \code{ERR\_NO\_TRANSACTION\textbackslash n} & Aucun chèque en cours \\
Caisse $\to$ lecteur & \code{ERR\_BUSY\textbackslash n} & Caisse occupée \\
\midrule
\multicolumn{2}{@{}l}{Trame vide, tronquée, corrompue ou absente} & Erreur à signaler \\
\bottomrule
\end{tabularx}
\end{table}

% =============================================================================================
\section{Démarche}

La démarche suit cinq étapes, chacune répondant à une question (figure~\ref{fig:demarche}).
L'étape~1 a mis en place le dépôt GitHub : une \emph{pull request} par étape, vérifiée par
l'intégration continue, et des commits qui justifient chaque changement.

\begin{figure}[H]
\centering
\begin{tikzpicture}[
    etape/.style={draw, rounded corners=2pt, minimum width=2.45cm, minimum height=0.9cm,
                  align=center, font=\small},
    question/.style={font=\footnotesize\itshape, align=center, text width=2.5cm},
    fleche/.style={-{Stealth[length=2.2mm]}, semithick},
    node distance=0.55cm,
  ]
  \node[etape] (g) {\textbf{2.} Générer};
  \node[etape, right=of g] (d) {\textbf{3.} Découpler};
  \node[etape, right=of d] (t) {\textbf{4.} Tester};
  \node[etape, right=of t] (m) {\textbf{5.} Mesurer};
  \node[etape, right=of m] (u) {\textbf{6.} Muter};
  \foreach \a/\b in {g/d, d/t, t/m, m/u} \draw[fleche] (\a) -- (\b);
  \node[question, below=3pt of g] {Le code existe-t-il ?};
  \node[question, below=3pt of d] {Est-il testable sans caisse ?};
  \node[question, below=3pt of t] {Fait-il ce qu'il faut ?};
  \node[question, below=3pt of m] {Tout est-il testé ?};
  \node[question, below=3pt of u] {Les tests sont-ils fiables ?};
\end{tikzpicture}
\caption{Les cinq étapes de la démarche et la question à laquelle chacune répond.}
\label{fig:demarche}
\end{figure}

% =============================================================================================
\section{Déroulement}

\subsection{Génération du module (étape 2)}

\paragraph{Démarche.} Un prompt précis
(\code{docs/prompts/etape-2-generation.md}) décrit le protocole, l'API attendue et les règles de
validation. Le code produit est intégré \emph{sans modification}.

\paragraph{Choix.} Bibliothèques imposées (\code{serialport}, \code{chrono}, \code{thiserror}),
montant en centimes entiers (\code{u64}) et une erreur explicite par cas (\code{PosError}).

\paragraph{Justification.} Un nombre flottant ne représente pas 145,50 exactement : pour une
somme d'argent, on compte en centimes.

\subsection{Découplage des dépendances (étape 3)}

\paragraph{Analyse.} Deux dépendances empêchent de tester le module
(tableau~\ref{tab:dependances}).

\begin{table}[H]
\centering
\caption{Dépendances fortes du code généré.}
\label{tab:dependances}
\begin{tabularx}{\linewidth}{@{}l l L@{}}
\toprule
Dépendance & Emplacement & Conséquence pour les tests \\
\midrule
Port série ouvert directement & Constructeur & Impossible de tester sans caisse \\
Horloge système & Validation de la date & Résultat variable selon le jour \\
\bottomrule
\end{tabularx}
\end{table}

\paragraph{Solution.} \emph{Injection de dépendances} : le module reçoit son port et son
horloge, \code{new(port, horloge)}, au lieu de les créer. Le port est un simple flux d'octets
(\code{Read + Write}), l'horloge une interface \code{Clock}. La fonction \code{open()} assemble
le vrai port et l'horloge système.

\paragraph{Justification.} Un faux port se résume à quelques octets, et la lecture de la trame
reste dans le code testé. Un exemple de documentation, exécuté par les tests, le démontre avec
un \emph{stub}.

\subsection{Tests unitaires (étape 4)}

\paragraph{Méthode.} Framework \code{mockall}. Les \emph{stubs} simulent la réponse de la
caisse et la date du jour ; les \emph{mocks} vérifient en plus les échanges (commande envoyée
une seule fois, aucune lecture si l'envoi échoue). Chaque cas est un test structuré en
\emph{Arrange / Act / Assert}.

\paragraph{Cas traités.} Cas nominaux ; erreurs de la caisse et de communication (timeout,
interruption) ; trames vides, tronquées ou corrompues ; dates spéciales (29~février, jour même,
lendemain, 31~avril) ; montants extrêmes (1~centime, valeur maximale, dépassement, zéro,
négatif).

\paragraph{Résultat.} Six tests échouent sur le code généré (phase rouge) ; trois corrections
les font passer (phase verte). Là où le protocole est muet, l'ambigu est refusé : \code{\textbackslash n}
final obligatoire, date strictement \code{AAAA-MM-JJ}, montants \code{145.} et \code{0145.50}
rejetés. Seule une date future est refusée.

\subsection{Couverture de code (étape 5)}

\paragraph{Méthode.} L'outil \code{cargo-llvm-cov} mesure les lignes et les
branchements\footnote{La mesure des branchements requiert la chaîne \emph{nightly} de Rust.}.
Chaque portion non couverte est soit une situation non testée, qui reçoit un test, soit du code
superflu, qui est supprimé ; un commit justifie chaque traitement
(tableau~\ref{tab:couverture}).

\begin{table}[H]
\centering
\caption{Portions non couvertes et traitement appliqué.}
\label{tab:couverture}
\begin{tabularx}{\linewidth}{@{}l l L@{}}
\toprule
Portion & Diagnostic & Traitement \\
\midrule
\code{SystemClock::today} & Non testée & Test ajouté \\
\code{open()}, échec et succès & Non testée & Port inexistant ; pseudo-terminal \\
Décimales non numériques & Branchement jamais pris & Cas \code{145.5a} \\
Partie entière hors \code{u64} & Non testée & Cas dédié \\
Erreur de calcul des centimes & Superflu (impossible) & Code simplifié \\
Programme \og hello world \fg{} & Superflu & Remplacé par une CLI \\
\bottomrule
\end{tabularx}
\end{table}

\paragraph{Résultat.} Lignes : 89\,\% $\to$ 100\,\% ; branchements : 97,5\,\% $\to$ 100\,\%.
Le fichier \code{src/main.rs}, code d'assemblage sans logique métier, est exclu de la mesure et
testé en lançant le programme.

\subsection{Tests de mutation (étape 6)}

\paragraph{Principe.} L'outil \code{cargo-mutants} introduit de petites erreurs dans le code
(\code{>} devient \code{>=}, \code{\&\&} devient \code{||}\dots). Un mutant qui survit révèle un
test trop faible.

\paragraph{Résultat.} Sur 63 mutants, 4 survivent ; chacun fait l'objet d'une issue, fermée par
le commit qui renforce les tests (tableau~\ref{tab:mutants}). Deux mutants provoquant un
\emph{timeout} sont aussi rendus détectables immédiatement. Campagne finale : aucun survivant.

\begin{table}[H]
\centering
\caption{Mutants survivants et renforcement des tests.}
\label{tab:mutants}
\begin{tabularx}{\linewidth}{@{}l l L@{}}
\toprule
Issue & Mutant survivant & Test ajouté \\
\midrule
\#11 & \code{\&\&} remplacé par \code{||} (format de date) & Date \code{2026-11-1} : seule la longueur est fausse \\
\#9, \#10 & \code{>} remplacé par \code{==} ou \code{>=} & Réponse de 256 octets exactement \\
\#8 & Programme principal vidé & Lancement de la CLI, codes de sortie vérifiés \\
\bottomrule
\end{tabularx}
\end{table}

% =============================================================================================
\section{Résultats}

\begin{table}[H]
\centering
\caption{Indicateurs avant et après traitement.}
\label{tab:resultats}
\begin{tabular}{@{}lrr@{}}
\toprule
Indicateur & Avant & Après \\
\midrule
Couverture des lignes & 89,3\,\% & 100\,\% \\
Couverture des branchements & 97,5\,\% & 100\,\% \\
Mutants survivants & 4 sur 63 & 0 sur 63 \\
Défauts du code généré & 6 & 0 \\
\bottomrule
\end{tabular}
\end{table}

La suite compte 60~tests : 55 pour le module, 1 pour l'horloge, 3 pour la CLI et 1 exemple de
documentation. Chaque \emph{pull request} est vérifiée par l'intégration continue : formatage,
analyse statique, tests sur Linux, Windows et macOS, version minimale de Rust, documentation,
couverture de 100\,\% des lignes et absence de mutant survivant.

% =============================================================================================
\section{Reproductibilité}

\noindent\begin{minipage}{\linewidth}
\begin{verbatim}
cargo test                                             # tests
cargo +nightly llvm-cov --branch \
      --ignore-filename-regex 'src/main\.rs' --open    # couverture
cargo mutants                                          # mutation
\end{verbatim}
\end{minipage}

% =============================================================================================
\section{Conclusion}

\begin{itemize}
  \item Le code généré semblait correct, mais acceptait six réponses invalides : seuls les
        tests l'ont révélé.
  \item L'injection de dépendances est ce qui rend les tests possibles, sans caisse réelle et
        à une date fixe.
  \item Une couverture de 100\,\% ne suffit pas : la mutation a encore trouvé quatre tests trop
        faibles, tous renforcés.
\end{itemize}

% =============================================================================================
\appendix
\section{Lexique}

\begin{description}[leftmargin=2.9cm, labelwidth=2.7cm, style=sameline, itemsep=1pt]
  \item[Test unitaire] Test d'un petit morceau de code isolé, rapide et au résultat constant.
  \item[AAA] \emph{Arrange, Act, Assert} : préparer, agir, vérifier.
  \item[Dépendance] Élément extérieur et imprévisible : port série, horloge, réseau.
  \item[Injection] Fournir ses dépendances au code au lieu qu'il les crée, pour les remplacer
        en test.
  \item[Stub, mock] Faux objet qui répond (\emph{stub}) ; qui vérifie en plus qu'on l'a bien
        appelé (\emph{mock}).
  \item[Couverture] Part du code exécutée par les tests.
  \item[Mutant] Copie du code avec une erreur volontaire, qui doit faire échouer un test.
  \item[TDD] Écrire le test d'abord (rouge), corriger le code (vert), puis remanier.
\end{description}

\end{document}
